\section{Introduction}

Large language models have recently demonstrated agentic capabilities---planning, tool invocation, and long-horizon reasoning---that make them attractive cores for intelligent assistants spanning chip design, verification, and other domains of electronic design automation.
However, state-of-the-art agent backbones are routinely 7\,B parameters or larger, demanding high-end accelerators and memory footprints that are incompatible with edge, embedded, and cost-sensitive deployments where inference latency and energy are first-class constraints.

Recent small language models (SLMs, $<\,$2\,B) have closed much of the gap with larger counterparts on single-turn benchmarks, but building such a model \emph{from scratch} remains an under-documented engineering problem.
Every choice---architecture, optimizer, learning-rate schedule, data mixture---interacts with every other, and a naive grid search is prohibitively expensive.
Our goal is therefore twofold: (i) to construct, from scratch and on a strictly limited compute budget, an end-to-end multimodal agentic system---an efficient hybrid language core paired with a custom vision encoder, aligned into a semiconductor-domain agent---and (ii) to expose the concrete, reproducible decisions that lead to it.

\paragraph{BaiZe.}
We present BaiZe, an end-to-end multimodal agentic system assembled from scratch across five stages: (i) pre-training a $\approx$2.2B-parameter hybrid (state-space + attention) LLM, (ii) post-training it for agentic capability, (iii) pre-training a custom vision encoder, (iv) aligning the two into a multimodal model, and (v) adapting the result to the semiconductor domain.
Three design pillars organise our work:

\begin{itemize}
\item \textbf{Architecture.} A Nemotron-H-style sequential hybrid language backbone---56 layers interleaving Mamba-2 state-space layers, four grouped-query-attention layers, and MLP layers---giving linear-time decoding for inference-heavy agents, paired with a from-scratch vision encoder.

\item \textbf{Reduced-horizon proxy search.} We exploit short, cheap proxy runs---in the spirit of maximal-update parameterization ($\mu$P) transfer~\cite{yang2022tensor}---so that the relative ranking of candidate hyper-parameters discovered at small scale carries over to the full model without re-tuning at scale.

\item \textbf{Budget-constrained stage-wise pipeline.} Each of the five stages is trained under a hard GPU-hour budget, with a staged search protocol (S1--S4) for the language backbone (numeric results reported as \textbf{[TBD]} pending re-measurement).
\end{itemize}

We hope BaiZe serves as a minimal, reproducible blueprint for building efficient agentic SLMs from scratch, and that our search methodology is of practical value to cost-bounded training efforts in the EDA and embedded-AI communities.